%%%%%%%%%%%%%%%%%%%%%%%%%%%%%%%%%%%%%%%%%%%%%%%%%%%%%%%%%%%%%%%%%%%%%%%%%%%
\chapter*{Conclusion}
\addcontentsline{toc}{chapter}{Conclusion}
%%%%%%%%%%%%%%%%%%%%%%%%%%%%%%%%%%%%%%%%%%%%%%%%%%%%%%%%%%%%%%%%%%%%%%%%%%%

Ce rapport est le reflet de notre stage de fin d'études effectué au sein de
\textbf{Yes Internet}, dans le cadre de l'obtention de notre licence en
Informatique de Gestion à l'\textbf{École Supérieure de l'Économie Numérique (ESEN)}.

Durant cette période, nous avons eu l'opportunité de mettre en pratique nos acquis
académiques à travers le développement d'une plateforme web de gestion de
location de véhicules. L'objectif était de proposer une solution centralisée,
fiable et évolutive, capable de faciliter la gestion des agences, des
clients, des véhicules et des réservations.

<<<<<<< HEAD
Pour repondre aux exigences du projet, nous avons suivi une demarche structurée.
Nous avons d'abord realise une analyse de l'existant, puis identifie les
besoins fonctionnels et non fonctionnels des differents acteurs. Ensuite,
nous avons adopte la methodologie SCRUM afin de planifier et executer le
travail en iterations: construction du backlog produit, planification des
sprints, implementation progressive et validation des fonctionnalites.

Sur le plan technique, notre solution s'appuie sur les outils principaux suivants:
Laravel, React, MySQL, Tailwind CSS et
Vite. Nous avons egalement utilise Git/GitHub pour le versioning,
Lucidchart pour la modelisation, Visual Studio Code pour le
developpement et XAMPP pour l'environnement local. Ces choix nous ont
permis de maintenir une architecture claire, une bonne qualite de code et un suivi
rigoureux du projet.

Ce stage a ete tres enrichissant, car il nous a permis de renforcer nos competences
techniques et organisationnelles, notamment en conception, en
developpement web full stack, en tests et en gestion du
temps. En perspective, la plateforme peut etre amelioree par l'integration de
fonctionnalites complementaires, telles que des notifications en temps
reel, des tableaux de bord plus avances et des mecanismes de
supervision securisee, afin d'optimiser davantage l'experience utilisateur
et la performance globale du systeme.
=======
Pour répondre aux exigences du projet, nous avons suivi une démarche structurée.
Nous avons d'abord réalisé une analyse de l'existant, puis identifié les
besoins fonctionnels et non fonctionnels des différents acteurs. Ensuite,
nous avons adopté la méthodologie SCRUM afin de planifier et exécuter le
travail en itérations : construction du backlog produit, planification des
sprints, implémentation progressive et validation des fonctionnalités.
>>>>>>> cc5484eb8c9f96bd8d3b034b8b0e3a775f34df87

Sur le plan technique, notre solution s'appuie sur les outils suivants :
\textbf{Laravel} pour le développement de l'API REST et la gestion de la
logique métier, \textbf{React} pour la construction des interfaces utilisateur
sous forme de SPA, \textbf{MySQL} pour le stockage et la gestion des données
relationnelles, \textbf{Tailwind CSS} pour la mise en forme et le design des
composants, et \textbf{Vite} pour la compilation et l'optimisation du frontend.
Nous avons également utilisé \textbf{Git/GitHub} pour le versioning du code,
\textbf{Lucidchart} pour la modélisation UML, \textbf{Visual Studio Code} comme
environnement de développement et \textbf{XAMPP} pour l'environnement local.
Ces choix nous ont permis de maintenir une architecture claire, une bonne
qualité de code et un suivi rigoureux du projet.

Ce stage a été une expérience professionnelle formatrice, qui nous a permis de consolider nos compétences en développement web full stack et de contribuer concrètement à la transformation numérique d'une agence tunisienne spécialisée dans les solutions web. Il nous a également préparés à intégrer le monde professionnel en nous confrontant à des contraintes réelles : respect des délais, communication avec le client, prise de décision technique et travail en équipe. Cette expérience constitue ainsi un tremplin solide vers une carrière dans le développement logiciel et l'ingénierie des systèmes d'information.

En perspective, plusieurs axes d'amélioration peuvent être envisagés pour
faire évoluer la plateforme. L'intégration d'un gateway de paiement en ligne
permettrait de digitaliser complètement le processus de réservation. Le suivi
du comportement des clients offrirait une meilleure compréhension des habitudes
d'utilisation et renforcerait le système de scoring de fiabilité. Enfin, le
développement d'une application mobile élargirait l'accessibilité de la
plateforme et améliorerait l'expérience utilisateur au quotidien.


